% Options for packages loaded elsewhere
% Options for packages loaded elsewhere
\PassOptionsToPackage{unicode}{hyperref}
\PassOptionsToPackage{hyphens}{url}
\PassOptionsToPackage{dvipsnames,svgnames,x11names}{xcolor}
%
\documentclass[
  brazilian,
  letterpaper,
  DIV=11,
  numbers=noendperiod]{scrreprt}
\usepackage{xcolor}
\usepackage{amsmath,amssymb}
\setcounter{secnumdepth}{5}
\usepackage{iftex}
\ifPDFTeX
  \usepackage[T1]{fontenc}
  \usepackage[utf8]{inputenc}
  \usepackage{textcomp} % provide euro and other symbols
\else % if luatex or xetex
  \usepackage{unicode-math} % this also loads fontspec
  \defaultfontfeatures{Scale=MatchLowercase}
  \defaultfontfeatures[\rmfamily]{Ligatures=TeX,Scale=1}
\fi
\usepackage{lmodern}
\ifPDFTeX\else
  % xetex/luatex font selection
\fi
% Use upquote if available, for straight quotes in verbatim environments
\IfFileExists{upquote.sty}{\usepackage{upquote}}{}
\IfFileExists{microtype.sty}{% use microtype if available
  \usepackage[]{microtype}
  \UseMicrotypeSet[protrusion]{basicmath} % disable protrusion for tt fonts
}{}
\makeatletter
\@ifundefined{KOMAClassName}{% if non-KOMA class
  \IfFileExists{parskip.sty}{%
    \usepackage{parskip}
  }{% else
    \setlength{\parindent}{0pt}
    \setlength{\parskip}{6pt plus 2pt minus 1pt}}
}{% if KOMA class
  \KOMAoptions{parskip=half}}
\makeatother
% Make \paragraph and \subparagraph free-standing
\makeatletter
\ifx\paragraph\undefined\else
  \let\oldparagraph\paragraph
  \renewcommand{\paragraph}{
    \@ifstar
      \xxxParagraphStar
      \xxxParagraphNoStar
  }
  \newcommand{\xxxParagraphStar}[1]{\oldparagraph*{#1}\mbox{}}
  \newcommand{\xxxParagraphNoStar}[1]{\oldparagraph{#1}\mbox{}}
\fi
\ifx\subparagraph\undefined\else
  \let\oldsubparagraph\subparagraph
  \renewcommand{\subparagraph}{
    \@ifstar
      \xxxSubParagraphStar
      \xxxSubParagraphNoStar
  }
  \newcommand{\xxxSubParagraphStar}[1]{\oldsubparagraph*{#1}\mbox{}}
  \newcommand{\xxxSubParagraphNoStar}[1]{\oldsubparagraph{#1}\mbox{}}
\fi
\makeatother


\usepackage{longtable,booktabs,array}
\usepackage{calc} % for calculating minipage widths
% Correct order of tables after \paragraph or \subparagraph
\usepackage{etoolbox}
\makeatletter
\patchcmd\longtable{\par}{\if@noskipsec\mbox{}\fi\par}{}{}
\makeatother
% Allow footnotes in longtable head/foot
\IfFileExists{footnotehyper.sty}{\usepackage{footnotehyper}}{\usepackage{footnote}}
\makesavenoteenv{longtable}
\usepackage{graphicx}
\makeatletter
\newsavebox\pandoc@box
\newcommand*\pandocbounded[1]{% scales image to fit in text height/width
  \sbox\pandoc@box{#1}%
  \Gscale@div\@tempa{\textheight}{\dimexpr\ht\pandoc@box+\dp\pandoc@box\relax}%
  \Gscale@div\@tempb{\linewidth}{\wd\pandoc@box}%
  \ifdim\@tempb\p@<\@tempa\p@\let\@tempa\@tempb\fi% select the smaller of both
  \ifdim\@tempa\p@<\p@\scalebox{\@tempa}{\usebox\pandoc@box}%
  \else\usebox{\pandoc@box}%
  \fi%
}
% Set default figure placement to htbp
\def\fps@figure{htbp}
\makeatother



\ifLuaTeX
\usepackage[bidi=basic]{babel}
\else
\usepackage[bidi=default]{babel}
\fi
% get rid of language-specific shorthands (see #6817):
\let\LanguageShortHands\languageshorthands
\def\languageshorthands#1{}


\setlength{\emergencystretch}{3em} % prevent overfull lines

\providecommand{\tightlist}{%
  \setlength{\itemsep}{0pt}\setlength{\parskip}{0pt}}



 


\KOMAoption{captions}{tableheading}
\makeatletter
\@ifpackageloaded{bookmark}{}{\usepackage{bookmark}}
\makeatother
\makeatletter
\@ifpackageloaded{caption}{}{\usepackage{caption}}
\AtBeginDocument{%
\ifdefined\contentsname
  \renewcommand*\contentsname{Índice}
\else
  \newcommand\contentsname{Índice}
\fi
\ifdefined\listfigurename
  \renewcommand*\listfigurename{Lista de Figuras}
\else
  \newcommand\listfigurename{Lista de Figuras}
\fi
\ifdefined\listtablename
  \renewcommand*\listtablename{Lista de Tabelas}
\else
  \newcommand\listtablename{Lista de Tabelas}
\fi
\ifdefined\figurename
  \renewcommand*\figurename{Figura}
\else
  \newcommand\figurename{Figura}
\fi
\ifdefined\tablename
  \renewcommand*\tablename{Tabela}
\else
  \newcommand\tablename{Tabela}
\fi
}
\@ifpackageloaded{float}{}{\usepackage{float}}
\floatstyle{ruled}
\@ifundefined{c@chapter}{\newfloat{codelisting}{h}{lop}}{\newfloat{codelisting}{h}{lop}[chapter]}
\floatname{codelisting}{Listagem}
\newcommand*\listoflistings{\listof{codelisting}{Lista de Listagens}}
\makeatother
\makeatletter
\makeatother
\makeatletter
\@ifpackageloaded{caption}{}{\usepackage{caption}}
\@ifpackageloaded{subcaption}{}{\usepackage{subcaption}}
\makeatother
\usepackage{bookmark}
\IfFileExists{xurl.sty}{\usepackage{xurl}}{} % add URL line breaks if available
\urlstyle{same}
\hypersetup{
  pdftitle={Análise de Dados para Biocientistas},
  pdfauthor={Márcio Souza},
  pdflang={pt-BR},
  colorlinks=true,
  linkcolor={blue},
  filecolor={Maroon},
  citecolor={Blue},
  urlcolor={Blue},
  pdfcreator={LaTeX via pandoc}}


\title{Análise de Dados para Biocientistas}
\author{Márcio Souza}
\date{2026-10-07}
\begin{document}
\maketitle

\renewcommand*\contentsname{Índice}
{
\hypersetup{linkcolor=}
\setcounter{tocdepth}{2}
\tableofcontents
}

\bookmarksetup{startatroot}

\chapter*{Bem-vindo}\label{bem-vindo}
\addcontentsline{toc}{chapter}{Bem-vindo}

\markboth{Bem-vindo}{Bem-vindo}

\textbf{Análise de Dados para Biocientistas} é um livro interativo
concebido para transformar a forma como os dados são analisados nas
ciências da vida.

\begin{quote}
\emph{Investir em equipamentos caros e robustos sem entender de dados é
como comprar um foguete para ir ao supermercado - parece impressionante,
mas no final, você ainda está perdido sem um plano de voo! A estatística
é a chave que dá acesso a novas rotas.}
\end{quote}

\subsection*{Ficha Técnica}\label{ficha-tuxe9cnica}
\addcontentsline{toc}{subsection}{Ficha Técnica}

\begin{itemize}
\tightlist
\item
  \textbf{Autor:} Márcio Luís Moreira de Souza
\item
  \textbf{ISBN:} 978-X-XXXX-XXXX-X
\item
  \textbf{DOI:}
  \href{https://doi.org/10.XXXX/zenodo.XXXXX}{10.XXXX/zenodo.XXXXX}
\item
  \textbf{Licença:} Este material é disponibilizado sob a licença
  Creative Commons Attribution-NonCommercial 4.0 International (CC BY-NC
  4.0).
\end{itemize}

\begin{center}\rule{0.5\linewidth}{0.5pt}\end{center}

Para iniciar a leitura e interagir com os códigos em tempo real, utilize
a barra de navegação lateral e avance para o \textbf{Capítulo 1:
Introdução ao método estatístico}.

\bookmarksetup{startatroot}

\chapter{Introdução ao método
estatístico}\label{introduuxe7uxe3o-ao-muxe9todo-estatuxedstico}

Este capítulo fornece uma introdução ao método estatístico{[}cite: 6{]}.
Ele discute os conceitos básicos de estatística, incluindo o que é o
método estatístico, os seus objetivos, tipos de dados, métodos de
recolha de dados e aplicações{[}cite: 6{]}.

\section{O que é o método
estatístico?}\label{o-que-uxe9-o-muxe9todo-estatuxedstico}

O método estatístico é um conjunto de procedimentos para recolher,
organizar, analisar e interpretar dados{[}cite: 7{]}. Ele é usado para
descrever e tirar conclusões sobre uma população a partir de uma
amostra{[}cite: 7{]}.

Os objetivos do método estatístico são:

\begin{itemize}
\tightlist
\item
  \textbf{Descrever dados:} pode ser usado para descrever os dados de
  uma população ou amostra, usando medidas de tendência central,
  dispersão e forma{[}cite: 7{]}.
\item
  \textbf{Explicar dados:} usado para identificar padrões e
  relacionamentos entre variáveis, usando métodos como regressão e
  análise de variância{[}cite: 7{]}.
\item
  \textbf{Prever dados:} usado para prever valores futuros, através de
  regressão e análise de séries temporais{[}cite: 7{]}.
\end{itemize}

O método estatístico permite-nos, a partir de um conjunto de símbolos
que representam quantidades ou medidas, transformá-los em dados para
extrair informações fundamentais para as atividades humanas{[}cite:
7{]}. \textbf{Dados} são um conjunto de números ou outras informações
que podem ser recolhidos, organizados e analisados{[}cite: 7{]}.
\textbf{Informação}, por sua vez, é um conjunto de dados que tem
significado e pode ser usado para tomar decisões{[}cite: 7{]}.

Dados podem ser estruturados ou não estruturados{[}cite: 7{]}. Dados
estruturados estão organizados num formato predefinido, como uma tabela
ou lista{[}cite: 7{]}. Dados não estruturados não possuem um formato
predefinido, como texto livre, imagens, áudio ou vídeo{[}cite: 7{]}. A
ciência de dados, a partir de métodos estatísticos, cria condições de
inovação e avanço do conhecimento humano{[}cite: 7{]}.

\section{Métodos de recolha de
dados}\label{muxe9todos-de-recolha-de-dados}

Os dados podem ser recolhidos de várias maneiras, incluindo:

\begin{itemize}
\tightlist
\item
  \textbf{Observação:} recolha de dados sem interferir na
  população{[}cite: 8{]}. Por exemplo, observação de doentes num
  hospital{[}cite: 8{]}.
\item
  \textbf{Experimentação:} recolha de dados intervindo na
  população{[}cite: 8{]}. Por exemplo, um ensaio clínico para testar a
  eficácia de um novo medicamento{[}cite: 8{]}.
\item
  \textbf{Levantamento:} recolha de dados por meio de questionários ou
  entrevistas{[}cite: 8{]}.
\end{itemize}

\section{Aplicações da
estatística}\label{aplicauxe7uxf5es-da-estatuxedstica}

A estatística tem uma ampla gama de aplicações:

\begin{itemize}
\tightlist
\item
  \textbf{Ciências da saúde:} usada para recolher, analisar e
  interpretar dados em estudos clínicos, epidemiológicos e de saúde
  pública{[}cite: 8{]}.
\item
  \textbf{Ciências sociais:} usada para recolher, analisar e interpretar
  dados em pesquisas sociais, económicas e políticas{[}cite: 8{]}.
\item
  \textbf{Ciências naturais:} usada para recolher, analisar e
  interpretar dados em pesquisas científicas, como física, química e
  biologia{[}cite: 8{]}.
\end{itemize}

\subsection{Exemplos nas ciências da
saúde}\label{exemplos-nas-ciuxeancias-da-sauxfade}

\begin{itemize}
\tightlist
\item
  Um estudo clínico realizado para testar a eficácia de um novo
  medicamento para tratar uma doença{[}cite: 8{]}.
\item
  Um estudo epidemiológico realizado para identificar fatores de risco
  para uma doença{[}cite: 8{]}.
\item
  Um estudo de saúde pública realizado para avaliar a eficácia de um
  programa de prevenção de doenças{[}cite: 8{]}.
\end{itemize}

\section{Tipos de dados (Variáveis)}\label{tipos-de-dados-variuxe1veis}

Os dados podem ser divididos em dois tipos principais:

\begin{itemize}
\tightlist
\item
  \textbf{Qualitativos:} dados que não podem ser medidos
  quantitativamente{[}cite: 9{]}. Por exemplo, dados sobre o sexo, a cor
  dos olhos ou a preferência política{[}cite: 9{]}.
\item
  \textbf{Quantitativos:} dados que podem ser medidos
  quantitativamente{[}cite: 9{]}. Por exemplo, dados sobre a altura, o
  peso ou a temperatura{[}cite: 9{]}.
\end{itemize}

Convencionalmente, em estatística, também chamamos aos dados
``variáveis'' em alusão à variedade de informação inerente à medida em
questão{[}cite: 9{]}. Logo, sexo, cor dos olhos, altura, peso, etc., são
exemplos de variáveis{[}cite: 9{]}.

O diagrama abaixo ilustra a classificação completa das variáveis de
forma dinâmica.

\includegraphics[width=12.08in,height=4.98in]{cap01_introducao_files/figure-latex/mermaid-figure-1.png}

Classificação das Variáveis Estatísticas

\bookmarksetup{startatroot}

\chapter{}\label{section}

\bookmarksetup{startatroot}

\chapter{}\label{section-1}

\bookmarksetup{startatroot}

\chapter{}\label{section-2}

\bookmarksetup{startatroot}

\chapter{}\label{section-3}

\bookmarksetup{startatroot}

\chapter{}\label{section-4}

\bookmarksetup{startatroot}

\chapter{}\label{section-5}

\bookmarksetup{startatroot}

\chapter{}\label{section-6}

\bookmarksetup{startatroot}

\chapter{}\label{section-7}

\bookmarksetup{startatroot}

\chapter{}\label{section-8}




\end{document}
